% =====================================================================
% CHAPTER 7: Design Methodology and Implementation (9-chapter template)
% Template source: thesis-or-project-template_1745300785.pdf, Chapter 7.
% Assembled per thesis/RESTRUCTURE_PLAN.md from:
%   - old mainmatter/ch3.tex: "Training Objective" (7.2.1, equations
%     byte-identical), "Datasets and Training Protocol" (7.2.2),
%     "Implementation and Deployment Environment" (into 7.3),
%     "Optimization of the Network for Edge Inference" (7.4).
%   - old appendices/app_a.tex "Software and Hardware Environment"
%     (into 7.3; Appendix A keeps the raw config table).
%   - old mainmatter/ch4.tex "Experimental Setup Summary" + the protocol
%     sentences of the measurement sections (7.6, 7.8; RESULTS stay in
%     Chapter 8).
% New prose: 7.1 Introduction, 7.5 Review and Feedback, 7.7 Principle of
% Operation, 7.9 Summary. No new numbers are introduced.
% Moved labels (each used once book-wide): eq:loss, fig:supervision,
% fig:input_protocol, tab:train_config, tab:efficiency_findings.
% Hard-coded "Chapter 4" refs in moved prose updated to Chapter 8.
% =====================================================================
\chapter{Design Methodology and Implementation}

\section{Introduction}
Chapter 3 fixed the architecture of StereoLite. This chapter covers how
that architecture became a trained, optimized, and measurable system: the
training objective and dataset protocol, the selection of
hyperparameters, hardware, and software, the optimization pass that met
the edge latency budget, the review and feedback process behind the
design, the experimental setup, the principle of operation from rectified
image pair to metric point cloud, and the measurement protocols behind
every number in Chapter 8. Chapter 8 reports and discusses the results;
this chapter covers only the methodology and implementation that produced
them.

\FloatBarrier
\section{Detailed Design Methodology}
The design methodology has two parts: the objective the network
minimizes during training, and the data and optimization protocol under
which that objective is applied.

\subsection{Training Objective}
The network trains under the multi-scale objective of \eqref{eq:loss}. A
pixel counts as valid when its ground-truth disparity is finite and lies
between 0 and 320 pixels. Each lower-scale prediction is resized to full
resolution and multiplied by its scale factor before comparison. The
objective weights an L1 term at every scale, heaviest at full
resolution, adds a matching L1 term on the geometry-volume output, and
augments the full-resolution prediction with a Sobel
gradient-consistency term, a threshold hinge stack, a KITTI-style
D1-relative hinge, an edge-aware smoothness term, and a gated
slant-supervision term that trains the plane slopes only where they are
trustworthy. The slant term pairs with plane rendering: it makes the
slopes accurate enough for the rendering step to sharpen boundaries.

The composition was chosen by ablation rather than guesswork. The loss
study of Chapter 8 tested nine formulations: plain L1 under-serves the
outlier metrics, hinge terms on a single threshold overfit that
threshold, and the threshold stack combined with the D1 hinge was the
only formulation to win both the end-point error and the outlier rates.
The production objective descends from that winner, keeping the
gradient term for its sub-pixel benefit and holding the smoothness term
to a small weight so it regularizes without blurring.

\begin{equation}
\begin{aligned}
\mathcal{L} \;=\;& 1.00\,\mathcal{L}_{1}(d_{\mathrm{full}})
 + 0.50\,\mathcal{L}_{1}(d_{1/2})
 + 0.30\,\mathcal{L}_{1}(d_{1/4})
 + 0.20\,\mathcal{L}_{1}(d_{1/8})
 + 0.10\,\mathcal{L}_{1}(d_{1/16})\\
&+ 0.15\,\mathcal{L}_{1}(d_{\mathrm{gev}})
 + 0.50\,\mathcal{L}_{\mathrm{grad}}
 + 0.20\,\mathcal{L}_{\mathrm{thr}}
 + 0.20\,\mathcal{L}_{\mathrm{D1}}
 + 0.02\,\mathcal{L}_{\mathrm{smooth}}
 + 0.30\,\mathcal{L}_{\mathrm{slant}}
\end{aligned}
\label{eq:loss}
\end{equation}
where $\mathcal{L}_1$ is the masked mean absolute disparity error,
$\mathcal{L}_{\mathrm{grad}}$ a Sobel gradient-consistency term,
$\mathcal{L}_{\mathrm{thr}}$ a hinge over the error thresholds
$\{0.5, 1, 2, 3\}$ pixels, $\mathcal{L}_{\mathrm{D1}}$ the D1-relative
hinge, $\mathcal{L}_{\mathrm{smooth}}$ an edge-aware smoothness term, and
$\mathcal{L}_{\mathrm{slant}}$ the gated slant supervision.

\begin{figure}[!htbp]
\centering
\includegraphics[width=\textwidth]{fig_3_8_supervision_preview}
\caption{\textbf{Multi-scale supervision.} Six disparity outputs receive weighted L1 terms against the ground truth per \eqref{eq:loss}; the full-resolution prediction additionally receives the gradient, threshold, D1, smoothness, and gated slant terms.}
\label{fig:supervision}
\end{figure}

\FloatBarrier
\subsection{Datasets and Training Protocol}
The model trains on the full SceneFlow finalpass dataset. The canonical
split has 35,454 training pairs, drawn from FlyingThings3D training
together with Monkaa and Driving; the FlyingThings3D test set of 4,370
pairs is held out for evaluation, since only FlyingThings3D provides an
official test split. Training samples random 384 by 640 windows,
co-located in both views, from the native 960 by 540 frames, which
preserves native disparity magnitudes. Chapter 8 shows this native-crop
protocol substantially outperforms the legacy protocol of downscaling
the whole frame. Three augmentations from the OpenStereo recipe apply:
asymmetric color jitter, right-image-only random erasing, and random
scaling. Optimization uses AdamW with a OneCycle schedule over 60,000
steps at a peak learning rate of $8 \times 10^{-4}$, a batch of 32, and
bfloat16 mixed precision, on a single A100 GPU. Input is normalized to
the unit interval at the model boundary. Validation uses a fixed
400-pair subset of the test set at the native 960 by 540 axis, evaluated
every thousand steps, with a resized-axis measurement logged alongside.
Table \ref{tab:train_config} lists the configuration.

\begin{figure}[!htbp]
\centering
\includegraphics[width=\textwidth]{fig_3_10_input_protocol}
\caption{\textbf{Native-crop input protocol.} Training samples random $384\times 640$ co-located windows of the native $960\times 540$ frame (left), preserving native disparity magnitudes; the legacy protocol (right) downscales the whole frame and shrinks every disparity by the same factor.}
\label{fig:input_protocol}
\end{figure}

\FloatBarrier
\section{Selection of Parameters, Apparatus and Components}

\begin{table}[H]
\centering
\caption{\textbf{Training configuration of the production run.} Values verbatim from the run's configuration record (meta.json of run 20260704).}
\label{tab:train_config}
\small
\begin{tabular}{ll}
\toprule
Element & Value\\
\midrule
Optimizer & AdamW, weight decay $10^{-5}$\\
Schedule & OneCycle, 60{,}000 steps, peak LR $8\times 10^{-4}$\\
Batch size / precision & 32 / bfloat16 mixed\\
Input protocol & random native $384\times 640$ co-located crops\\
Augmentation & asymmetric color jitter, right-only erase, random scale\\
Training data & full SceneFlow finalpass, 35{,}454 pairs\\
Validation & fixed 400-pair FlyingThings3D subset, every 1k steps\\
Hardware & one NVIDIA A100 80 GB (cloud)\\
\bottomrule
\end{tabular}
\end{table}

Each hyperparameter is the result of a controlled study reported in
Chapter 8, and the design-evidence audit of Chapter 3 traces every choice
to its justification. The principal architectural parameters are the
truncated YOLO26s backbone kept through its first seven blocks, a
32-channel recurrent hidden state, a 48-channel refinement head, eight
correlation groups over 24 disparity hypotheses at one sixteenth
resolution, a local correlation half-range of two pixels, a refinement
schedule of two, three, and three iterations across the three scales, and
a geometry encoding volume narrowed to a band of plus or minus 16 around
the tile disparity. The training-side parameters appear in Table
\ref{tab:train_config}. Table \ref{tab:arch_config} of Appendix A gives
the complete configuration record of the trained model, verbatim from its
checkpoint.

The apparatus divides into compute and sensing. The model is implemented
in PyTorch 2.11 with CUDA 12.8 on Python 3.10. Training runs on a single
cloud A100 GPU with 80 gigabytes of memory. Ablations and evaluations run
on a cloud T4, chosen for its low cost, which supports a controlled study
after every design change. The local latency benchmark runs on an RTX
3050 laptop GPU with roughly 4 gigabytes of memory, representing the low
end of the deployment memory budget. The intended deployment target is
the Jetson Orin Nano class of edge accelerator, on which the optimized
model was measured directly; Chapter 8 reports the on-device result. On
the sensing side, real-camera pairs come from a calibrated low-cost
binocular rig that streams a 2560 by 720 side-by-side image, split into a
rectified 1280 by 720 pair per frame. This rig supplies the real-data
qualification of the trained model, including the zero-shot qualitative
pairs, the real-camera agreement protocol of Section
\ref{sec:testing_protocols}, and the three-dimensional reconstructions of
Chapter 8.

\FloatBarrier
\section{Optimization and Configuration Selection}
Reaching the latency target required optimizing the network's inference
graph. Six changes were applied: rewriting the cost-volume construction
loops to use padded zero-copy views instead of allocating and zeroing a
feature map for every disparity bin; reducing the grid-sample calls from
48 to eight by batching the lookups and caching the sampling grids;
hoisting the static channels of the recurrent input, which are
re-convolved every iteration, once per scale, following the
precomputed-context trick of RAFT-Stereo \cite{lipson2021raftstereo};
fusing the gate convolutions and the four output heads, which share
inputs; removing dead inference-path code; and narrowing the geometry
volume to a band around the tile disparity. Each change was proven
equivalent before adoption, either bitwise or through a matched A/B
comparison, and together they cut the full-precision latency by a factor
of 1.74. The pass also cleared three operators that block eight-bit
deployment: the unfold in convex upsampling, group normalization on some
accelerator toolchains, and three-dimensional convolution on neural
processing units. Each was swapped for a quantization-friendly
equivalent, and the resulting graph was exported and calibrated to an
eight-bit TensorRT engine for the Jetson Orin Nano, whose measured
on-device latency Chapter 8 reports. Table \ref{tab:efficiency_findings}
summarizes the changes.

\begin{table}[!htbp]
\centering
\caption{\textbf{Inference optimizations and their equivalence status.} Savings measured at $384\times 640$ on RTX 3050; equivalence class states how each change was proven safe before adoption.}
\label{tab:efficiency_findings}
\small
\begin{tabular}{p{5.6cm}p{4.2cm}l}
\toprule
Optimization & Saving & Equivalence\\
\midrule
Zero-copy padded views in cost-volume loops & 5 to 15\% latency & bitwise\\
Batched correlation lookups, cached grids & 3 to 8\% latency & identical math\\
Static GRU input channels hoisted per scale & 8 to 10\% of MACs & exact\\
Fused gate convolutions and output heads & 1 to 3\% latency & bitwise\\
Dead inference-path code removed & memory + launches & exact\\
Narrow GEV band ($\pm 16$ around tile $d$) & 25 to 30\% of MACs & matched A/B (Ch.~8)\\
\midrule
Combined & fp32 61.4 ms vs 106.7 ms base ($1.74\times$) & \\
\bottomrule
\end{tabular}
\end{table}

\FloatBarrier
\section{Review and Feedback}
The implementation followed a review-and-feedback process rather than
unchecked iteration: every proposed change, architectural, loss-term, or
input-protocol, had to win a controlled ablation on the matched harness
of Chapter 3 before entering the production configuration, and a change
that lost its comparison was not adopted no matter how plausible it
looked. Every run, win or lose, went into a master experiments log with
its full configuration and metrics, so the evidence behind each
decision is auditable and no comparison depends on memory. Training
curves, qualitative panels, and ablation tables were reviewed with the
supervisor at regular intervals, and that feedback redirected the work
more than once, most visibly toward the cross-dataset evaluation gate
that now follows every training run. Failed variants were documented
rather than discarded: the rejected mechanisms and anti-synergistic
compositions of Chapter 8 constrain future work as firmly as the
adopted ones justify the present design.

\FloatBarrier
\section{Simulation and Experimental Setup}
This section describes the setup behind every experiment in Chapter 8,
which reports the results. Three datasets are used. SceneFlow finalpass
supplies the synthetic training corpus and in-domain test bed under the
split and protocol of Table \ref{tab:train_config}: 35,454 training
pairs and a held-out FlyingThings3D test set of 4,370 pairs, with a
fixed 400-pair subset used for validation and the perturbation sweep.
Middlebury 2014, 23 perfect-set scenes of real indoor imagery the model
never trains on, is the zero-shot cross-domain test. The third source is
the project's physical stereo rig: the calibrated 2560 by 720
side-by-side camera of Section 7.3, which supplies real rectified pairs
(no ground truth) for qualitative evaluation, teacher agreement, and
three-dimensional reconstruction. Figure \ref{fig:setup} summarizes the
setup: these data sources, the training and evaluation platforms, and
the recorded metrics.

\begin{figure}[H]
\centering
\includegraphics[width=\textwidth]{fig_7_setup_preview}
\caption{\textbf{Simulation and experimental setup.} Three data sources (synthetic SceneFlow with dense ground truth, the real zero-shot Middlebury 2014 benchmark, and the physical stereo rig) feed the StereoLite network. The model is trained on a cloud A100 and evaluated on a cloud T4 for accuracy, on an RTX 3050 laptop for latency, and on the Jetson Orin Nano deployment target for on-device latency. Every driver logs metrics and metadata that Chapter 8 reports.}
\label{fig:setup}
\end{figure}

Each experiment runs as a separate scripted driver on a cloud GPU: one
for full-test-set evaluation, one for zero-shot Middlebury 2014
evaluation (using the same protocol for the proposed model and the published
reference checkpoints), and one for the rectification-robustness sweep.
Every driver writes its configuration and metrics to a metadata file
alongside its outputs, and Appendix A names the specific scripts so each
reported number can be reproduced. A cross-domain evaluation run on a
cloud T4 costs only minutes and a few cents. That low cost keeps the
after-every-training-run evaluation gate of Section 7.5 practical.
Accuracy is measured everywhere with the standard stereo metrics of
\eqref{eq:metrics}, averaged per image and then over images; the
protocols for latency, rectification robustness, and real-camera
agreement are given in Section \ref{sec:testing_protocols}. On-device
latency is measured on the Jetson Orin Nano developer board of Figure
\ref{fig:jetson}, the network's intended edge deployment target.

\begin{figure}[H]
\centering
\includegraphics[width=0.6\textwidth]{fig_7_jetson}
\caption{\textbf{Jetson Orin Nano developer board.} The edge accelerator used as the on-device deployment target, on which the optimized eight-bit model is timed (Chapter 8). It provides a few tera-operations per second within a power budget of a few watts to a few tens of watts.}
\label{fig:jetson}
\end{figure}

\FloatBarrier
\section{Principle of Operation}
In deployment the system operates as a fixed feed-forward pipeline: a
rectified left and right pair enters at the inference resolution of
384 by 640 pixels, normalized to the unit interval, and passes through
the blocks of Chapter 3. The shared
encoder processes both views in a single batched pass and emits matching
features at one quarter, one eighth, and one sixteenth resolution, while
a separate context encoder reads the left image alone and produces the
context maps that feed the recurrent state.

Matching begins at the coarsest scale. The tile initialization module
builds the group-wise correlation volume of \eqref{eq:gwc} at one
sixteenth resolution, regularizes it with a small three-dimensional
network, and reads off the initial disparity and confidence through the
soft-argmax of \eqref{eq:softargmax_init}; the plane slopes start at
zero, and the hidden state starts from the coarsest context map. The
recurrent refinement runs its fixed schedule: two iterations at one
sixteenth, three at one eighth, and three at one quarter resolution. At
each iteration the right features are warped by the current disparity, a
local correlation is evaluated over five offsets, and the convolutional
gated recurrent unit of \eqref{eq:convgru} updates the hidden state;
gated residuals adjust the disparity, the plane slopes, and the
confidence. Between scales, the plane-aware propagation of
\eqref{eq:plane_prop} carries the tile state to the finer grid along its
local slanted-surface hypothesis rather than by plain interpolation.

At one quarter resolution, the narrow-band geometry encoding volume
produces its independent proposal per \eqref{eq:gev_expect}, and the
fail-soft gate of \eqref{eq:failsoft} blends that proposal into the
recurrent estimate only where the learned weight favors it. This
corrects regions the tile state resolved poorly, without disturbing the
ones it resolved well. The refined quarter-resolution state is rendered
from its tile planes with an edge gate and upsampled to full
resolution by the two learned convex-combination stages, each output
subpixel a softmax-weighted mixture of its three-by-three coarse
neighborhood. The result is a dense full-resolution disparity map.

Disparity becomes depth through the calibrated relations of
\eqref{eq:backprojection}: depth follows $Z = fB/d$, and each valid pixel
optionally back-projects to a metric three-dimensional point in the
camera frame, colored by its left-image pixel. On the project's rig, the
calibrated focal length and baseline scale to the inference resolution by
the width ratio, so the point cloud is metric without further tuning.
This single forward pass, from image pair to disparity map and optional
colored point cloud, is the operation measured throughout Chapter 8.

\FloatBarrier
\section{Testing and Measurement}
\label{sec:testing_protocols}
Chapter 8 reports numbers produced under four measurement protocols,
defined here so that each has a stated procedure. The numbers themselves
appear only in that chapter.

The first protocol measures latency and memory. Latency is measured at
batch one and 384 by 640 resolution after warm-up passes; memory is
recorded as the peak inference working set. Every measurement states its
device. On the RTX 3050 laptop GPU the model runs in eager PyTorch at
full and half precision, without compiler assistance, so the numbers
reflect the unoptimized serving path. On the deployment target, the
Jetson Orin Nano, the optimized model is exported to an eight-bit
TensorRT engine and measured directly on the device; this is the figure
of merit for the real-time claim.

The second protocol tests tolerance to imperfect rectification under
controlled conditions. The trained model is evaluated at the native axis
on the fixed 400-pair FlyingThings3D subset. The right image is shifted
vertically by a controlled offset while the disparity ground truth stays
unchanged, so the offset emulates a residual vertical misalignment
between the views. The sweep runs from zero through sub-pixel offsets to
several pixels, and the standard metrics are reported at each point.

The third protocol corroborates the synthetic sweep on physical data.
The model runs on 997 rectified pairs from the project's stereo rig,
whose factory rectification leaves a residual vertical misalignment, and
its disparity is compared against the FoundationStereo teacher's
pseudo-ground-truth on the same pairs. Because both the model and the
teacher see the same misaligned input, this protocol measures agreement
with a strong reference rather than absolute error; it corroborates,
rather than replaces, the controlled sweep.

The fourth is the dual-axis evaluation convention. Every in-domain
evaluation is reported on the native 960 by 540 axis, which preserves
the true disparity magnitudes; a resized-axis measurement at 640 by 384
is logged alongside for comparability with the legacy protocol.
Input-protocol comparisons evaluate every arm on both axes, so a gain on
one axis cannot hide a loss on the other. Zero-shot Middlebury 2014
evaluation follows a single fixed protocol applied the same way to the
model and to the public reference checkpoints: inference at 384 by 640,
ground truth scaled to the inference width, disparities above 192 pixels
masked, and per-scene metrics are macro-averaged over scenes.

\FloatBarrier
\section{Summary}
This chapter traced the path from the architecture of Chapter 3 to a
trained, deployable system. The network trains under a multi-scale
composite objective chosen by ablation, using a native-crop SceneFlow
protocol, the three OpenStereo augmentations, a OneCycle schedule, and a
cloud A100. Every hyperparameter, apparatus choice, and software
component rests on documented evidence; the full configuration appears
in Appendix A. An equivalence-proven optimization pass brought the
inference graph within the edge latency budget. Controlled ablations
gated every design decision, a master experiments record captured each
run, and the supervisor reviewed the work regularly. The experimental
setup covers a synthetic training and test corpus, an unseen real
benchmark, and the physical stereo rig. Four measurement protocols
(latency, rectification perturbation, real-camera agreement, and
dual-axis evaluation) define how Chapter 8's numbers were obtained.
